\documentclass[12pt]{article}
\usepackage[margin=1in]{geometry}
\usepackage{amsmath, amssymb, amsthm}

\title{\textbf{Advanced Probability Theory}\\
GU 4156: Spring 2025\\
\large Homework \#2\\
\small Exercises 5, 6(i), 7, 8 are due by 10 PM Friday, Feb. 28}
\author{Professor Ioannis Karatzas\\
Department of Mathematics, Columbia University}
\date{February 17, 2025}

\begin{document}
\maketitle

\noindent In Exercises 1--3, all statements are with respect to a given measurable space and a given, fixed filtration. In Exercises 4--9 we introduce also a probability measure.

\bigskip

\noindent\textbf{Exercise \#1:} If $S$ and $T$ are stopping times, then
\begin{enumerate}
    \item[(i)] $S \vee T$ and $S \wedge T$ are stopping times;
    \item[(ii)] $S + T$ is a stopping time.
\end{enumerate}

\bigskip

\noindent\textbf{Exercise \#2:} For any stopping times $S$ and $T$, we have
\begin{enumerate}
    \item[(i)] $\mathcal{F}_S \subseteq \mathcal{F}_T$, if $S \leq T$;
    \item[(ii)] $\mathcal{F}_{S \wedge T} = \mathcal{F}_S \cap \mathcal{F}_T$.
\end{enumerate}

\bigskip

\noindent\textbf{Exercise \#3:} Let $X = \{X_0, X_1, \cdots\}$ be an adapted sequence of random variables, and $T : \Omega \to \mathbb{N}_0$ a (finite) stopping time. Define the random variable $X_T$, and argue that it is measurable with respect to $\mathcal{F}_T$.

\bigskip

\noindent\textbf{Exercise \#4: A Simple but Useful Optional Sampling Result.} On a probability space $(\Omega, \mathcal{F}, P)$ endowed with a filtration $\mathbb{F} = (\mathcal{F}_n)_{n \in \mathbb{N}_0}$, consider a martingale $M = (M_n)_{n \in \mathbb{N}_0}$ and a stopping time $T$ with $P(T < \infty) = 1$.
\begin{enumerate}
    \item[(a)] Show that $E\left[M_T\right] = E\left[M_0\right]$ holds then under the additional conditions
    \[
        E\left[|M_T|\right] < \infty, \quad \lim_{n \to \infty} E\left[M_n \mathbf{1}_{\{T > n\}}\right] = 0.
    \]
    \item[(b)] Show by example that, in general,$^1$ $E\left[M_T\right] = E\left[M_0\right]$ need not hold.
\end{enumerate}

\bigskip

\noindent\textbf{Exercise \#5: Yet Another Simple Optional Sampling Result.} Suppose that $T$ is a stopping time with $E(T) < \infty$, and that $X = \{X_0, X_1, \cdots\}$ is a martingale (resp. super, sub) for which, with some $K \in (0, \infty)$,
\[
    E\left[X_{n+1} - X_n \,\middle|\, \mathcal{F}_n\right] \leq K \quad \text{holds a.e. on } \{T \geq n\}, \; \forall\, n \in \mathbb{N}_0.
\]
Then we have $E(|X_T|) < \infty$ and $E(X_T) = E(X_0)$ (resp. $\leq$, $\geq$).

\smallskip
\noindent(\textit{Hint}: Argue first, that $E(X_{T \wedge m}) = E(X_0)$ (resp. $\leq$, $\geq$) holds for every $m \in \mathbb{N}_0$; then find an upper bound for $E\left[|X_T - X_{T \wedge m}|\right]$ that goes to zero as $m \to \infty$.)

\bigskip

\noindent\textbf{Exercise \#6: The Identities of A. Wald.} Suppose that $\xi_1, \xi_2, \cdots$ are independent random variables, of common distribution with $E(|\xi_1|) < \infty$; and that the underlying filtration is the one generated by this sequence. Fix also a stopping time $\tau$ with $E(\tau) < \infty$.
\begin{enumerate}
    \item[(i)] Show that $E\left[\sum_{n=1}^{\tau} \xi_n\right] = E(\tau) \cdot E(\xi_1)$.
    \item[(ii)$^*$] (Optional) If, in addition, $E(\xi_1^2) < \infty$, show that
    \[
        E\left[\left(\sum_{n=1}^{\tau} \xi_n - \tau \cdot E(\xi_1)\right)^2\right] = E(\tau) \cdot \mathrm{Var}(\xi_1).
    \]
\end{enumerate}

\bigskip

\noindent\textbf{Exercise \#7: Switching between Supermartingales.} Let $X = \{X_n\}_{n \in \mathbb{N}_0}$ and $Y = \{Y_n\}_{n \in \mathbb{N}_0}$ be supermartingales, and $\tau$ a stopping time such that $X_\tau \geq Y_\tau$ holds a.e.\ on $\{\tau < \infty\}$. Show that the sequence $Z = \{Z_n\}_{n \in \mathbb{N}_0}$ defined by
\[
    Z_n := X_n \cdot \mathbf{1}_{\{\tau > n\}} + Y_n \cdot \mathbf{1}_{\{\tau \leq n\}}, \quad n \in \mathbb{N}_0
\]
is again a supermartingale.

\bigskip

\noindent\textbf{Exercise \#8: The Maximal Inequality of J. Ville for Positive Supermartingales.} For every nonnegative supermartingale $X = \{X_n\}_{n \in \mathbb{N}_0}$, the random variable $\sup_{n \in \mathbb{N}_0} X_n$ is a.e.\ finite, and satisfies for $\lambda > 0$ the inequality
\[
    P\left(\sup_{n \in \mathbb{N}_0} X_n > \lambda \;\middle|\; \mathcal{F}_0\right) \leq 1 \wedge \frac{X_0}{\lambda}.
\]
(\textit{Hint}: Consider the first time the level $\lambda$ is exceeded, and use Exercise 7.)

\bigskip

\noindent\textbf{Exercise \#9: Simple Random Walk with Negative Drift.} Consider the simple random walk $S = (S_n)_{n \in \mathbb{N}_0}$ on the integers $\mathbb{Z}$, started at some fixed site $S_0 \equiv a \in \mathbb{Z}$, with probability $p \in (0, 1/2)$ of moving up and probability $1-p$ of moving down on every step. We are denoting by $S_n$ the position of the walk on day $t = n$.$^2$

\begin{enumerate}
    \item[(a)] Show $S = (S_n)_{n \in \mathbb{N}_0}$ is a Markov Chain; specify the transition probabilities.

    \item[(b)] Show that $S_\infty := \lim_{n \to \infty} S_n = -\infty$ holds a.e.; and that the ``maximum-height'' random variable $Q := \sup_{n \in \mathbb{N}_0} S_n$ is a.e.\ finite: $P(Q < \infty) = 1$.

    \item[(c)] Show that the random sequence $M = (M_n)_{n \in \mathbb{N}_0}$ given by
    \[
        M_n = \left(\frac{1-p}{p}\right)^{S_n}, \quad n = 0, 1, 2, \cdots \tag{0.1}
    \]
    is a martingale, and specify the filtration for which this statement is true.

    \item[(d)] For some given integers $N \geq 2$ and $a \in \{1, 2, \cdots, N-1\}$, let us assume that the walk starts at the site $S_0 = a$, and consider the first time
    \[
        T := \min\{n \geq 1 : S_n = 0 \text{ or } N\} \tag{0.2}
    \]
    the walk reaches one of the sites $0$ or $N$ (always with the understanding that $T = \infty$, if this never happens). Show that with $r := (1-p)/p$, we have then
    \[
        P(S_T = 0) = \frac{r^a - r^N}{1 - r^N}, \quad E(S_T) = N \cdot \frac{1 - r^a}{1 - r^N}.
    \]
    (\textit{Hint}: Recall the result of Exercise 4(a), and the martingale of part (c) here.)

    \smallskip
    \noindent For the remainder of this Exercise, let us start the random walk at $S_0 \equiv 0$ for simplicity.

    \item[(e)] Use the Strong Markov property, to show that the maximum-height $Q$ in part (b) here has geometric distribution
    \[
        P(Q \geq k) = \vartheta^k, \quad k \in \mathbb{N} \tag{0.3}
    \]
    with parameter $\vartheta := P(Q > 0)$.

    \item[(f)] With $\tau := \min\{n \geq 1 : S_n = 1\}$, argue that $(M_{n \wedge \tau})_{n \in \mathbb{N}_0}$ is a bounded martingale. Conclude from this, that the parameter of the geometric distribution in (e) is
    \[
        \vartheta = \frac{p}{1-p}.
    \]

    \item[(g)] What is the probability that this random walk, started at the origin, never rises above it to become positive?
\end{enumerate}

\vspace{1cm}
\noindent\rule{0.5\textwidth}{0.4pt}

\noindent $^1$ I.e., without additional conditions.

\noindent $^2$ It is always a good idea to draw pictures of such things.

\end{document}
